\documentclass[12pt, a4paper]{article}

% ============================================================
% BROAD LITERATURE REVIEW
% Project: Constructing Hedge Portfolios with Neural Networks
% Purpose: Broad, critical literature review (3 000--4 000 words)
%
% TEAM STRATEGY
% - Assign one lead per major section.
% - Keep one editor responsible for tone, notation, transitions, deduplication.
% - Do not allow each subsection to become a disconnected paper summary.
%
% HOW TO WRITE EACH SUBSECTION
% 1. Define the theme or branch of the literature.
% 2. Explain what problem the branch is solving.
% 3. Compare the key papers.
% 4. State strengths, assumptions, and weaknesses.
% 5. Explain what remains unresolved.
%
% DO NOT WRITE LIKE THIS:
%   "Paper A did X. Paper B did Y. Paper C did Z."
%
% INSTEAD WRITE LIKE THIS:
%   "Recent work falls into two main camps... The first... The second...
%    The key difference is... The main limitation is..."
% ============================================================

% ------------------------------------------------------------
% PACKAGES
% ------------------------------------------------------------
\usepackage[T1]{fontenc}
\usepackage[utf8]{inputenc}
\usepackage[margin=2.5cm]{geometry}
\usepackage{amsmath, amssymb}
\usepackage{hyperref}
\usepackage{natbib}
\usepackage{microtype}
\usepackage{setspace}
\usepackage{parskip}
\usepackage{authblk}
\usepackage{xcolor}
\usepackage{booktabs}
\usepackage{tabularx}
\usepackage{longtable}
\usepackage{array}
\usepackage{enumitem}
\usepackage{mdframed}   % for the assignment boxes

\hypersetup{
    colorlinks=true,
    linkcolor=blue!70!black,
    citecolor=green!50!black,
    urlcolor=blue!70!black,
}

\setstretch{1.12}

% ---- Coloured assignment box ----
\newmdenv[
    backgroundcolor=yellow!15,
    linecolor=orange!60,
    linewidth=1pt,
    innertopmargin=4pt,
    innerbottommargin=4pt,
    innerleftmargin=6pt,
    innerrightmargin=6pt,
    skipabove=4pt,
    skipbelow=4pt,
]{assignbox}

% ---- Utility macros ----
\newcommand{\TODO}[1]{\textcolor{red}{[TODO: #1]}}
\newcommand{\paperlink}[2]{\href{#2}{#1}}
\newcommand{\crit}[1]{\noindent\textbf{Critical point.} #1}
\newcommand{\litquestion}[1]{\noindent\textbf{Question to answer.} \emph{#1}}

% ============================================================
% TITLE
% ============================================================
\title{%
    {\large Master of Financial Engineering: Research}\\[0.5em]
    \textbf{Broad Literature Review: Deep Learning for Derivative Hedging}%
}

\author[1]{Micaela TODO-SURNAME}
\author[1]{Kaiden TODO-SURNAME}
\author[1]{Siphelele TODO-SURNAME}
\author[1]{Glasson Osborne}
\affil[1]{African Institute of Financial Markets and Risk Management}

\date{\today}

% ============================================================
\begin{document}
% ============================================================

\maketitle

% ============================================================
% WORD COUNT TARGET: 3 000 -- 4 000 WORDS
% Run: texcount lit_review.tex   (or use Overleaf's word count)
% ============================================================

% ============================================================
% INTERNAL GUIDE  (delete before submission if desired)
% ============================================================

\section*{Internal guide for the group}

This document is designed to help produce an \emph{excellent} literature review rather
than a merely adequate one. The central aim is to explain, evaluate, and organise the
literature on neural networks for derivative hedging in a way that makes the research
gap and the logic of the project obvious.

\textbf{Core writing principle:} the review must be \emph{critical and taxonomic}. That
means the literature should be organised by themes, conceptual differences, modelling
assumptions, and practical implications --- not by publication date.

\textbf{Core questions this review should answer:}
\begin{enumerate}[label=\arabic*.]
    \item Why is hedging difficult in practice, even though classical theory provides
          elegant benchmarks such as Black--Scholes delta hedging?
    \item How did the literature move from pricing networks to direct hedging networks?
    \item What exactly is meant by \emph{deep hedging}, and why has it become the
          central modern framework?
    \item Under what conditions do neural approaches appear to outperform classical
          benchmarks?
    \item Are those improvements economically meaningful and robust, or are they
          sometimes driven by stylised simulation choices?
    \item What unresolved question or gap motivates the present project?
\end{enumerate}

\textbf{Recommended reading priority:}
\begin{enumerate}[label=\arabic*.]
    \item \citet{RufWang2020}
    \item \citet{BuehlerEtAl2019}
    \item \citet{HutchinsonLoPoggio1994}
    \item \citet{GarciaGencay2000}
    \item \citet{RufWang2022}
    \item \citet{DuEtAl2020}
    \item \citet{HorvathTeichmannZuric2021}
    \item \citet{Cont2025}
\end{enumerate}

% ============================================================
% DIVISION OF LABOUR TABLE  (keep or delete)
% ============================================================

\section*{Division of labour}
\begin{table}[h!]
\centering
\begin{tabularx}{\textwidth}{@{}l l X@{}}
\toprule
Section & Lead author & Notes / papers \\
\midrule
Introduction and motivation    & \TODO{Name} & Set scope; explain why deep hedging matters \\
Classical foundations          & \TODO{Name} & Black--Scholes, incomplete markets, quadratic hedging \\
Pricing vs hedging networks    & \TODO{Name} & HLP 1994, Garcia--Gen\c{c}ay 2000, conceptual transition \\
Deep hedging core              & \TODO{Name} & Buehler et al.\ 2019, Ruf--Wang 2022 \\
Extensions and synthesis       & \TODO{Name} & DRL, rough vol, data-driven hedging, critique \\
Final editor                   & \TODO{Name} & Tone, notation, transitions, deduplication \\
\bottomrule
\end{tabularx}
\end{table}

\tableofcontents
\newpage

% ============================================================
% 1. INTRODUCTION AND SCOPE
% ============================================================
\section{Introduction and Scope}
\label{sec:intro}

\begin{assignbox}
\textbf{Assignee:} \underline{\hspace{4cm}} \quad Target: $\approx 350$ words
\end{assignbox}

% TODO:
% - Define the scope of the literature review.
% - State that the emphasis is on hedging.
% - Explain that the review is broad but organised around a taxonomy.
% - Signal that the final section will identify the gap motivating the project.

Deep learning has become increasingly prominent in mathematical finance, including
applications to derivative pricing, calibration, risk management, and hedging. Among
these applications, hedging is especially important because it lies at the intersection
of theory and practice. A hedging rule must perform not merely in an idealised model,
but under discrete trading, model uncertainty, frictions, and often incomplete market
structure. This creates a natural environment in which flexible function approximators
such as neural networks may be useful.

This literature review focuses on neural-network methods for derivative hedging, with
particular emphasis on the modern \emph{deep hedging} paradigm. Rather than presenting
the literature chronologically, we organise it into conceptual branches. We begin with
the classical benchmark of complete-market replication and the practical reasons why
imperfect hedging arises. We then distinguish early neural-network methods for pricing
and sensitivity estimation from later approaches that learn hedge policies directly.
The central section of the review evaluates the deep hedging framework introduced by
\citet{BuehlerEtAl2019}, after which we discuss recent extensions involving
reinforcement learning, rough volatility, richer architectures, frictions, and
data-driven market generation. The review concludes with a comparative synthesis and a
statement of the research gap relevant to the present project.

The goal is not simply to catalogue papers, but to assess what the literature has
actually established. In particular, a strong review should ask whether neural methods
offer genuine economic advantages over classical hedging benchmarks, under what
assumptions these gains appear, and whether the evidence is robust to more realistic
data and benchmark choices. This critical perspective is essential because a literature
review should justify the project's direction, not merely provide background.

\litquestion{What makes hedging a more demanding application of machine learning than
pricing alone?}

% ============================================================
% 2. CLASSICAL FOUNDATIONS: WHY HEDGING IS DIFFICULT IN PRACTICE
% ============================================================
\section{Classical Foundations: Why Hedging Is Difficult in Practice}
\label{sec:foundations}

\begin{assignbox}
\textbf{Assignee:} \underline{\hspace{4cm}} \quad Target: $\approx 600$ words
\end{assignbox}

% TODO:
% - Explain Black--Scholes delta hedging as the benchmark.
% - Explain why perfect replication breaks down in practice.
% - Briefly classify imperfect-hedging approaches (taxonomy, not chronology).
% - Motivate why approximation methods become attractive.

\subsection{Complete-Market Replication as the Benchmark}
\label{subsec:bsm}

Classical derivative pricing theory provides the natural starting point for any
discussion of hedging. In the Black--Scholes framework \citep{BlackScholes1973,
Merton1973}, a contingent claim can be replicated by dynamic trading in the underlying
asset and a riskless account, yielding both a unique no-arbitrage price and an
associated delta-hedging strategy. This benchmark remains central in the literature
not because it is fully realistic, but because it provides the cleanest reference
point against which more flexible or data-driven hedging methods are evaluated.

% TODO: Explain the relationship between no-arbitrage pricing, replication, and delta
% hedging. Keep the focus on why this benchmark matters for later neural-network work.

\subsection{Why Perfect Hedging Fails Outside the Idealised Benchmark}
\label{subsec:imperfect}

The assumptions behind exact replication are restrictive. In practice, hedging is
discrete rather than continuous, volatility is not constant, jumps and stochastic
volatility may matter \citep{Heston1993}, market frictions such as transaction costs
cannot be ignored \citep{Leland1985}, and markets may be incomplete. In such settings,
there is generally no exact replicating strategy, and the hedger must choose how to
trade off replication accuracy against other objectives such as variance, downside
risk, or transaction cost control.

% TODO: Quantify the hedging error introduced by discrete rebalancing.
% Cite: Bertsimas, Kogan \& Lo (2000).

\subsection{Major Classical Approaches to Imperfect Hedging}
\label{subsec:classical_approaches}

Rather than listing papers chronologically, the classical literature organises into a
taxonomy of objectives:
\begin{enumerate}[label=\alph*.]
    \item \textbf{Super-replication:} eliminate short-position risk by ensuring the
          portfolio dominates the payoff in all admissible scenarios, at the cost of
          conservatism and often very high prices.
    \item \textbf{Utility or indifference-based hedging:} optimise a utility or
          certainty-equivalent criterion, making the hedger's preferences explicit but
          also creating dependence on an often hard-to-specify utility function.
    \item \textbf{Quadratic or mean-variance hedging:} minimise a squared
          replication-error criterion, obtaining a mathematically tractable framework
          that is especially relevant for later deep hedging work.
          \citet{Stangroom2021} provides a useful departmental example of neural
          approximation in this quadratic-hedging setting, benchmarking against
          dynamic programming, Black--Scholes, and Heston solutions.
\end{enumerate}

% TODO: For each approach, answer: What is the objective? What risk notion is
% optimised? Is the strategy self-financing? What are the strengths and practical
% difficulties?

\subsection{Why This Classical Background Matters}
\label{subsec:classical_motivation}

Once hedging is formulated as a sequential optimisation problem under imperfect market
conditions, approximation methods become attractive. Neural networks then enter the
literature not as replacements for classical finance, but as flexible tools for
approximating pricing rules, sensitivities, or trading policies when the classical
setting is no longer analytically tractable.

\crit{This section should be concise. Its purpose is to motivate modern neural
hedging, not to become a full essay on all of mathematical finance.}

% ============================================================
% 3. FROM PRICING NETWORKS TO HEDGING NETWORKS
% ============================================================
\section{From Pricing Networks to Hedging Networks}
\label{sec:pricing_to_hedging}

\begin{assignbox}
\textbf{Assignee:} \underline{\hspace{4cm}} \quad Target: $\approx 550$ words
\end{assignbox}

% TODO:
% - Explain the early pricing-network literature.
% - Distinguish it from direct hedging.
% - Prepare the reader for deep hedging.

\subsection{Early Neural Networks for Derivative Pricing}
\label{subsec:early_pricing}

The earliest neural-network literature in this area largely treated derivative pricing
as a nonlinear regression problem. Under this paradigm, the network received option
and market variables as inputs and produced an option price as output. Hedging, if
considered at all, was usually obtained indirectly by differentiating the learned
pricing function with respect to the underlying state variables \citep{HutchinsonLoPoggio1994}.

This pricing-first paradigm is historically important because it introduced
nonparametric function approximation into derivative markets, but it differs
fundamentally from modern deep hedging. In particular, the object learned is a pricing
map rather than a trading policy.

% TODO: Discuss Hutchinson, Lo \& Poggio (1994) and Garcia \& Gen\c{c}ay (2000)
% in depth. Garcia--Gen\c{c}ay introduced a homogeneity hint, showing that
% domain knowledge can improve network performance.
% Cite Ruf \& Wang (2020) as the main survey source classifying these contributions.

\subsection{The Conceptual Weakness of Pricing-First Hedging}
\label{subsec:pricing_limits}

A central critical point must be made here: accurate pricing does not automatically
imply accurate hedging. A network may fit option prices well while still producing
unstable or economically poor hedge ratios when differentiated. Small errors in a
pricing surface can translate into materially different sensitivities, and two models
with similar pricing performance may imply meaningfully different trading strategies.

This observation is especially important in hedging problems, where the action taken
by the trader is the hedge itself, not merely a price estimate. It therefore motivates
the shift from learning prices to learning hedge policies directly.

% TODO: Support with evidence from Ruf \& Wang (2020). Make the transition to
% direct policy learning explicit.

\subsection{Direct Hedge-Output Approaches as the Bridge}
\label{subsec:direct_output}

Some contributions between the early pricing literature and modern deep hedging began
to move toward networks that output sensitivities or hedge ratios directly. This is a
useful bridge in the literature because it shows that the field gradually recognised
that the hedge should be treated as a distinct prediction target rather than a
by-product of a pricing model.

\crit{The transition sentence at the end of this section should make it clear that
deep hedging represents the next conceptual step: instead of learning a price or a
Greek, the network learns a sequential trading policy.}

% ============================================================
% 4. DEEP HEDGING AS THE CENTRAL MODERN FRAMEWORK
% ============================================================
\section{Deep Hedging as the Central Modern Framework}
\label{sec:deep_hedging}

\begin{assignbox}
\textbf{Assignee:} \underline{\hspace{4cm}} \quad Target: $\approx 800$ words
\end{assignbox}

% TODO:
% - Explain Buehler et al. (2019) carefully.
% - Describe what deep hedging learns.
% - Explain why it is attractive.
% - Critique it rather than merely praising it.

\subsection{What Deep Hedging Is}
\label{subsec:dh_core}

The modern deep hedging literature is centred on the idea that hedging should be
treated as a sequential decision problem. In this framework, a neural network
parameterises the hedge policy itself. Given a market state at each trading time, the
network outputs the hedge position, and the parameters are trained by optimising a
chosen risk functional of terminal profit and loss over simulated or generated market
paths \citep{BuehlerEtAl2019}.

This formulation differs fundamentally from early pricing-network approaches because
the learned object is the trading rule. The framework also accommodates frictions,
constraints, and alternative objectives more naturally than the classical analytical
benchmark.

% TODO: Self-financing portfolio, discounted P\&L, recurrent NN parameterisation.
% When trained on BS paths with MSE loss, the network approximately recovers the
% BS delta. Discuss the role of N (time steps) and M (paths).

\subsection{Why the Framework Is Attractive}
\label{subsec:dh_attractive}

% TODO: Expand each point into a sentence or two.
\begin{itemize}
    \item It can incorporate transaction costs and constraints directly into the
          objective.
    \item It does not require a closed-form pricing formula.
    \item It is naturally applicable to multi-asset, path-dependent, or otherwise
          analytically difficult settings.
    \item It allows training directly against economically meaningful loss or risk
          criteria.
\end{itemize}

\subsection{What the Literature Claims}
\label{subsec:dh_claims}

The literature commonly claims that deep hedging can match or outperform classical
hedging rules, especially once frictions or model misspecification are present. It is
also often presented as a flexible approximation framework for settings in which
classical delta hedging is no longer optimal or even well-defined.

% TODO: Summarise the main empirical claims without yet endorsing them uncritically.
% Extensions: Carbonneau (2021) for long-dated FX options; Horvath et al. (2021)
% for rough volatility; Stangroom (2021) for mean-variance objective with RNN,
% TCN, AttentionNet and SpanMLP architectures -- RNN achieved best loss and VaR.

\subsection{What Must Be Evaluated Critically}
\label{subsec:dh_critical}

A strong literature review must ask several critical questions:
\begin{itemize}
    \item Is the benchmark fair? If the market truly follows Black--Scholes without
          frictions, analytical delta hedging is already a very strong baseline.
    \item Are the gains largest only in settings where frictions, constraints, or
          misspecification are introduced?
    \item Are the reported gains economically large, or merely statistically
          detectable?
    \item How dependent are the results on architecture, hyperparameter tuning,
          and feature engineering? \citet{Stangroom2021} provides direct evidence
          of architecture sensitivity: comparing RNN, TCN, AttentionNet, and SpanMLP
          under entropic loss and mean-P\&L objectives, the RNN achieves the best
          performance while remaining competitive even with shorter input spans.
    \item Does simulation-based training generalise to real market data?
    \item Is the learned strategy interpretable or stable enough for practical use?
\end{itemize}

\subsection{A Key Comparator: Simpler Models Versus Neural Networks}
\label{subsec:dh_comparator}

\citet{RufWang2022} is especially valuable because it tempers strong claims of neural
superiority. One of the major lessons from this literature is that neural networks
should not be treated as automatically dominant. In some settings, simpler
regression-based methods with sensible features remain highly competitive. This forces
benchmark discipline and helps distinguish genuine advances from over-parameterised
enthusiasm.

\crit{This section is the heart of the review. It should explain what deep hedging
solves, what it assumes, and when it appears genuinely more compelling than classical
hedging benchmarks.}

% ============================================================
% 5. RECENT EXTENSIONS AND FRONTIER DIRECTIONS
% ============================================================
\section{Recent Extensions and Frontier Directions}
\label{sec:extensions}

\begin{assignbox}
\textbf{Assignee:} \underline{\hspace{4cm}} \quad Target: $\approx 800$ words
\end{assignbox}

% TODO: Organise by themes, NOT by publication date.
% Compare branches rather than listing papers.
% Show where the frontier of the field seems to be moving.

\subsection{Richer Dynamics and Non-Markovian Environments}
\label{subsec:rough_vol}

One important direction is the extension of deep hedging beyond simple Markovian
Black--Scholes-type worlds. \citet{HorvathTeichmannZuric2021} examine settings
involving rough volatility or richer path structure. These contributions matter because
they test whether neural architectures can exploit temporal information and
non-Markovian dependence in ways that closed-form benchmarks cannot capture easily.

% TODO: Explain why such settings are a genuinely harder test for neural hedgers
% than the standard Black--Scholes case.

\subsection{Reinforcement Learning Formulations}
\label{subsec:rl}

A closely related branch frames the problem more explicitly in reinforcement-learning
language. Here the market state is treated as the state variable, the hedge as the
action, and the hedging objective as the reward or return criterion \citep{DuEtAl2020}.

A good literature review should clarify the relationship between this branch and the
broader deep hedging literature. In some papers the distinction is mainly one of
language and training setup, while in others it reflects a more explicit
control-theoretic or RL formulation.

\subsection{Frictions, Constraints, and Alternative Objectives}
\label{subsec:frictions}

For practical finance, this branch may be the most important. The strongest case for
deep hedging arguably emerges not when analytical replication is already available,
but when real-world features such as transaction costs, position constraints, market
impact, or alternative risk criteria are introduced. In such settings, the flexibility
of a learned policy becomes more valuable.

% TODO: Do not merely state that frictions can be included. Explain WHY this is one
% of the principal economic motivations for the field.
%
% Stangroom (2021) is directly relevant here:
%   - Uses entropic risk measure as loss function (not just MSE).
%   - Also studies mean P&L as an objective -- minimising expected profit and loss.
%   - Compares RNN, TCN, AttentionNet and SpanMLP under these objectives.
%   - RNN achieves best VaR values and lowest entropic loss.
%   - Longer input spans improve all architectures; RNN remains competitive even
%     with shorter spans.
% Use this to illustrate how loss function and input structure interact with
% architecture choice.

\subsection{Data-Driven and Model-Free Hedging}
\label{subsec:data_driven}

Another frontier concerns the data on which neural hedgers are trained. Classical
simulation-based work often trains on paths generated by a specified model such as
Black--Scholes or Heston. This is convenient, but it can also build the answer into
the experiment. Real markets provide only one realised history, creating a serious
data bottleneck.

This motivates interest in richer market generators and data-driven hedging
frameworks, such as \citet{Cont2025}. This branch is especially relevant because it
addresses one of the most persistent weaknesses of simulation-based finance papers:
impressive results under synthetic data may not translate into genuine out-of-sample
value.

\subsection{The Stangroom Dissertation as a Departmental Synthesis Source}
\label{subsec:stangroom}

The UCT dissertation by \citet{Stangroom2021} is useful as a \emph{supporting
synthesis source}, especially on deep hedging under a quadratic or mean-variance
objective, benchmark design against dynamic programming and Heston solutions, and the
use of pseudo-real or generated market paths. It should not replace the original
literature, but it can help interpret how the field is operationalised in a local
dissertation context.

\crit{A high-quality review should compare these branches and ask where the field is
really advancing: better architectures, more realistic objectives, stronger benchmarks,
or more realistic data.}

% ============================================================
% 6. COMPARATIVE SYNTHESIS
% ============================================================
\section{Comparative Synthesis}
\label{sec:synthesis}

\begin{assignbox}
\textbf{Assignee:} \underline{\hspace{4cm}} \quad Target: $\approx 500$ words
\end{assignbox}

% TODO: DO NOT introduce many new papers here. Synthesise the literature across
% common dimensions. Show that the group has understood the field as a structured
% body of work.

\subsection{Useful Comparison Dimensions}
\label{subsec:comparison}

A clear synthesis can compare papers across the following axes:
\begin{itemize}
    \item \textbf{Objective:} mean squared hedging error, convex risk measure, daily
          mark-to-market P\&L, tail-focused objective.
    \item \textbf{Training data:} Black--Scholes simulation, Heston simulation,
          rough-volatility paths, generated pseudo-real data, real market data.
    \item \textbf{Output type:} option price, implied quantity, Greek, direct hedge
          ratio, whole sequential policy.
    \item \textbf{Benchmark:} Black--Scholes delta, Heston hedge, linear model,
          dynamic-programming benchmark, no-hedge baseline.
    \item \textbf{Frictions:} absent, transaction costs only, richer constraints
          and impact.
    \item \textbf{Architecture:} feedforward, semi-recurrent, recurrent,
          weight-sharing, path-aware models.
    \item \textbf{Economic significance:} whether improvements are materially
          meaningful or only marginal.
\end{itemize}

\subsection{What the Synthesis Shows}
\label{subsec:synthesis_findings}

% TODO: Turn the prompts below into polished prose paragraphs.

\begin{quote}
The literature suggests that the principal contribution of deep hedging is not that
neural networks replace classical finance, but that they provide a flexible
optimisation framework once the assumptions of continuous, frictionless, correctly
specified replication are relaxed.
\end{quote}

\begin{quote}
A recurring weakness in the literature is that gains are often demonstrated under
stylised simulations, making it difficult to distinguish true economic advantage from
favourable experimental design.
\end{quote}

\begin{quote}
Recent work increasingly shifts attention away from raw predictive flexibility and
toward benchmark fairness, realistic frictions, richer objectives, and data realism.
\end{quote}

\crit{This section should sound like interpretation, not note-taking.}

% ============================================================
% 7. RESEARCH GAP AND RELEVANCE TO THE PRESENT PROJECT
% ============================================================
\section{Research Gap and Relevance to the Present Project}
\label{sec:gap}

\begin{assignbox}
\textbf{Assignee:} \underline{\hspace{4cm}} \quad Target: $\approx 400$ words
\end{assignbox}

% TODO:
% 1. Summarise what the literature has already established.
% 2. State what remains unclear, contested, or weakly evidenced.
% 3. Explain exactly what the present project will investigate.
%
% Possible gap statements:
%   - whether neural hedging still adds value when benchmark assumptions are
%     approximately correct;
%   - whether reported gains persist under discrete trading with fair comparison
%     to analytical deltas;
%   - whether alternative architectures (recurrent networks, weight sharing)
%     improve generalisation;
%   - whether deeper models are especially valuable under Heston / rough volatility
%     rather than Black--Scholes;
%   - whether the choice of loss function materially changes the learned hedge.

Despite the rapid progress surveyed above, several important problems remain open
\ldots

% Template paragraph -- replace with your own:
%
% Overall, the literature indicates that neural-network hedging is most promising
% when hedging is treated as a direct sequential decision problem under imperfect
% market conditions. However, uncertainty remains regarding the extent to which
% reported improvements reflect genuine economic advantages rather than favourable
% modelling assumptions, stylised simulations, or weak benchmark design. The present
% project contributes to this literature by \TODO{state your exact project angle
% here}, thereby assessing \TODO{state the main question here} in a setting that is
% directly comparable to the most relevant existing work.

% ============================================================
% INTERNAL APPENDIX: CORE PAPERS AND LINKS
% (delete before submission if desired)
% ============================================================

\section*{Internal appendix: core papers and links}

\subsection*{Priority reading list}
\begin{longtable}{>{\raggedright\arraybackslash}p{0.30\textwidth}
                  >{\raggedright\arraybackslash}p{0.64\textwidth}}
\toprule
Paper & Link(s) \\
\midrule
\endfirsthead
\toprule
Paper & Link(s) \\
\midrule
\endhead
\citet{RufWang2020}, \emph{Neural networks for option pricing and hedging: a literature review} &
\paperlink{arXiv}{https://arxiv.org/abs/1911.05620} \\[4pt]

\citet{BuehlerEtAl2019}, \emph{Deep Hedging} &
\paperlink{arXiv}{https://arxiv.org/abs/1802.03042};
\paperlink{Journal page}{https://www.tandfonline.com/doi/abs/10.1080/14697688.2019.1571683} \\[4pt]

\citet{HutchinsonLoPoggio1994}, \emph{A Nonparametric Approach to Pricing and Hedging Derivative Securities via Learning Networks} &
\paperlink{NBER working paper}{https://www.nber.org/papers/w4718};
\paperlink{Wiley page}{https://onlinelibrary.wiley.com/doi/abs/10.1111/j.1540-6261.1994.tb00081.x} \\[4pt]

\citet{GarciaGencay2000}, \emph{Pricing and Hedging Derivative Securities with Neural Networks and a Homogeneity Hint} &
\paperlink{Journal page}{https://www.sciencedirect.com/science/article/abs/pii/S0304407699000184};
\paperlink{Accessible PDF}{https://myrenegarcia.com/wp-content/uploads/2017/04/jecon2000.pdf} \\[4pt]

\citet{RufWang2022}, \emph{Hedging with Linear Regressions and Neural Networks} &
\paperlink{arXiv}{https://arxiv.org/abs/2004.08891};
\paperlink{Journal page}{https://www.tandfonline.com/doi/abs/10.1080/07350015.2021.1931241} \\[4pt]

\citet{DuEtAl2020}, \emph{Deep Reinforcement Learning for Option Replication and Hedging} &
\paperlink{PDF}{https://cims.nyu.edu/~ritter/du2020deep.pdf};
\paperlink{Journal page}{https://www.pm-research.com/content/iijjfds/2/4/44} \\[4pt]

\citet{HorvathTeichmannZuric2021}, \emph{Deep Hedging under Rough Volatility} &
\paperlink{arXiv}{https://arxiv.org/abs/2102.01962};
\paperlink{MDPI page}{https://www.mdpi.com/2227-9091/9/7/138} \\[4pt]

\citet{Cont2025}, \emph{Data-driven hedging with generative models} &
\paperlink{Springer page}{https://link.springer.com/article/10.1007/s10479-025-06867-3} \\[4pt]

\citet{Stangroom2021}, \emph{Deep Hedging in Incomplete Markets} (UCT MSc dissertation) &
\paperlink{arXiv}{https://arxiv.org/abs/2112.10084} \\
\bottomrule
\end{longtable}

\subsection*{How to use the Stangroom dissertation}

Use \citet{Stangroom2021} selectively as a support source for:
\begin{itemize}
    \item mean-variance or quadratic hedging context,
    \item benchmark design against dynamic programming, Black--Scholes, or Heston hedges,
    \item architecture comparison (RNN, TCN, AttentionNet, SpanMLP) under entropic
          loss and mean-P\&L objectives,
    \item discussion of pseudo-real or generated market paths,
    \item examples of how a departmental dissertation structures the story from
          classical theory into deep hedging.
\end{itemize}

Do \textbf{not} let it replace the original literature. Treat it as a secondary
synthesis source rather than a primary anchor citation.

% ============================================================
% INTERNAL APPENDIX: ONE-PAGE NOTE TEMPLATE
% (copy for each paper while reading; delete before submission)
% ============================================================

\section*{Internal appendix: one-page note template for each paper}

\begin{itemize}
    \item \textbf{Citation:}
    \item \textbf{One-sentence summary:}
    \item \textbf{Problem setting:}
    \item \textbf{Model learns:} price / Greek / hedge ratio / sequential policy
    \item \textbf{Training data:} simulated / generated / real
    \item \textbf{Objective:}
    \item \textbf{Benchmark:}
    \item \textbf{Main contribution:}
    \item \textbf{Main limitation:}
    \item \textbf{Why relevant to our project:}
    \item \textbf{Which section of the literature review it belongs in:}
\end{itemize}

% ============================================================
% FINAL CHECKLIST
% ============================================================

\section*{Final checklist before submission}
\begin{itemize}
    \item Is the review broad enough for 3\,000--4\,000 words but still clearly
          focused on hedging?
    \item Is the structure taxonomic rather than chronological?
    \item Does each section contain comparison, evaluation, and limitations?
    \item Is the difference between pricing networks and direct hedging networks made
          explicit?
    \item Is Deep Hedging treated as the core modern framework?
    \item Do we discuss benchmark fairness, frictions, objective choice, and data
          realism?
    \item Does the final section identify a precise research gap tied to our project?
    \item Has one editor removed duplication and harmonised notation and tone?
\end{itemize}

% ============================================================
% REFERENCES
% ============================================================
\newpage
\bibliographystyle{plainnat}
\bibliography{../paper/references}   % shared BibTeX file

\end{document}
